\chapterimage{street_plaque.jpg}

\chapter{On ets?}\label{sec:where}

Connectar dos ordinadors els fa molt més útils que dues màquines aïllades. 
Tanmateix, el seu veritable potencial només es desbloqueja quan es poden connectar a \textit{moltíssims} altres ordinadors.

Limitacions pràctiques com el cost i la ubicació geogràfica fan que sigui impossible connectar
directament cada parell d'ordinadors. En lloc d'això, els dispositius es distribueixen en un graf de xarxes interconnectades,
que es pot estendre fins i tot més enllà de l'escala d'un planeta.

Tant si és localment, a l'abast físic, com en un altre continent, cal un sistema d'adreçament
per identificar, localitzar i encaminar els missatges perquè arribin als peers (parells) previstos, tal com
fem amb les ubicacions físiques.



\section{Si us plau, contacta'm}\label{sec:where:topology}

Quan només hi intervenen $2$ \conceptRef{device}{dispositius}, es pot fer servir una connexió \conceptRef{point-to-point}{punt a punt}. En aquest tipus de connexió,
només hi ha ``aquest costat'' i ``l'altre costat'' del cable. A més, quan un dispositiu vol comunicar-se amb l'altre extrem, 
només hi ha un enllaç (\eg, un cable) on triar, de manera que no hi ha cap confusió possible. Com que podem fer funcionar aquesta connexió 
amb només $2$ \conceptRef{network card}{targetes de xarxa} i $1$ cable, les connexions punt a punt són viables per a $N=2$ dispositius.

\begin{wrapfigure}{R}{0.25\linewidth}
\includegraphics[width=\linewidth]{complete_graph.pdf}
\end{wrapfigure}


Si volem connectar $N > 2$ dispositius, una opció ingènua és connectar cada dispositiu amb tots els altres
en un \conceptRef{complete graph}{\textit{graf complet}}. Ara cada dispositiu ha de gestionar $N - 1$ cables cap als altres dispositius, 
però comunicar-se amb qualsevol altre dispositiu és tan fàcil com triar el cable correcte i establir una connexió punt a punt.

El problema principal d'aquest enfocament és que el nombre total de connexions és $N \cdot (N-1) / 2 = O(N^2)$
(en saber més al teu estimat curs de \textit{Matemàtica Discreta}). Això vol dir que, per tenir $N=10$ dispositius connectats d'aquesta manera,
necessitaries per exemple $40$ cables \textit{i} $90$ \conceptRef{network card}{targetes de xarxa}. Per a $N=1000$ dispositius, 
hi hauria mig milió de connexions punt a punt, cosa que ja és força ingestionable.

\begin{remark}
$O(N^2)$ és la notació d'una fita de complexitat \conceptRef{quadratic}{quadràtica}. 
\textit{No} és el mateix que \conceptRef{exponential}{exponencial} i no s'han de confondre mai.
% 
\readMore{https://en.wikipedia.org/wiki/Big_O_notation}
\end{remark}

\begin{wrapfigure}{L}{0.17\linewidth}
\includegraphics[height=4cm]{bus_graph.pdf}
\end{wrapfigure}

Tenim un truc amagat a la màniga: els \conceptRef{bus}{busos de dades}, en què una o més línies de dades
estan connectades simultàniament a diversos dispositius. 
Això conserva la simplicitat de les connexions punt a punt, perquè cada dispositiu només ha de gestionar
$1$ cable i el pot fer servir per comunicar-se amb tots els altres dispositius. També es conserva 
la rendibilitat, ja que només calen $N$ targetes de xarxa per a $N$ dispositius ($O(N)$). 

El principal avantatge dels busos de dades és també la seva principal debilitat: 
tots els dispositius envien i reben dades per la mateixa línia de dades, 
però no ho poden fer alhora perquè només hi ha una línia. Si dos o més ho fan, es produeix
una \conceptRef{collision}{col·lisió} que fa impossible la transmissió de dades mentre dura.\\[-0.5cm]

\begin{remark}
A mesura que es desenvolupaven les xarxes d'ordinadors, es van provar moltes alternatives a la 
\conceptRef{topology}{topologia} de bus. Un exemple curiós és \concept{Token Ring}, en què els dispositius
estan connectats al voltant d'un cercle i els missatges hi circulen en un sol sentit.
\end{remark}

A la pràctica, no és trivial implementar una tecnologia de bus 
on es puguin afegir nous dispositius amb el temps. 
% 
Històricament, de vegades es feia perforant físicament el cable principal del bus
i afegint un cable nou cap a l'ordinador que s'hi connectava. 
Més endavant, això es va simplificar amb uns dispositius anomenats \conceptRef{hub}{hubs} (concentradors). 
Funcionaven de manera semblant, però oferien \conceptRef{port}{ports} ja integrats 
i connectors extraïbles.
% 
Avui dia, els hubs han estat substituïts pels \conceptRef{switch}{switches} (commutadors). 
Els switches ofereixen la funcionalitat d'un bus, connectant cada dispositiu amb tots els altres,
i a més eviten les col·lisions. Per raons òbvies, estructurar les connexions 
al voltant d'un switch s'anomena \conceptRef{star topology}{topologia en estrella}.

\begin{exercise}
Què \textbf{cal en un switch} que no calgui en un hub?

\begin{center}
\includegraphics[width=0.75\linewidth]{switch.jpg}
\end{center}
\end{exercise}

Les topologies de bus funcionen molt bé a escala local i s'utilitzen àmpliament en
\conceptRef{LAN}{\conceptRef{network}{xarxes} d'àrea local (LANs)}. Tanmateix, no són rendibles
(de vegades ni tan sols físicament possibles) quan els dispositius són lluny els uns dels altres.
% 
En lloc d'això, els dispositius es poden organitzar en LANs separades, cadascuna conceptualment un bus, 
i després aquestes LANs separades es poden connectar mitjançant connexions punt a punt
més barates i factibles. Considerats en conjunt, aquests dispositius formarien llavors una 
\conceptRef{WAN}{xarxa d'àrea extensa (WAN)}.

\begin{center}
\includegraphics[width=0.65\linewidth]{lan_wan_structure.pdf}
\end{center}


\section{Necessito trobar-te}\label{sec:where:addresses}

Un cop més, el principal avantatge dels \conceptRef{bus}{busos} de dades és també la seva principal debilitat: 
tots els dispositius envien i reben dades per la mateixa línia de dades. 
Encara que no hi hagi col·lisions, cada missatge enviat al bus el 
reben tots els dispositius.
Tanmateix, el que realment volem és que el nostre missatge arribi al seu destí, 
i només al seu destí. Anomenem-ho el \textit{problema de trobar el meu dispositiu}.

El problema de trobar el meu dispositiu no és exclusiu dels busos o de les \conceptRef{LAN}{LANs}. Al contrari, esdevé 
molt important quan considerem les \conceptRef{WAN}{WANs}, on no tots els dispositius estan interconnectats directament.
En aquest cas, hem d'assegurar-nos que podem \conceptRef{routing}{encaminar} els nostres missatges
entre diferents LANs i que arriben al seu destí.

Sorprenentment, ni tan sols les connexions punt a punt estan a resguard del problema de trobar el meu dispositiu.
Suposant que els dispositius A i B estan connectats d'aquesta manera, dos programes diferents poden voler 
intercanviar informació al mateix temps, \eg, un programari de control d'alarmes i 
una aplicació de reproducció de música en temps real. 
En aquest escenari, vols que el \concept{client} i el \conceptRef{server}{servidor} del \conceptRef{service}{servei} 
de monitoratge de l'alarma es comuniquin entre ells, però no amb el servei de música en temps real, 
i viceversa.

La solució més habitual a aquest problema és fer servir identificadors que ens permetin
diferenciar entre dispositius o fins i tot entre \conceptRef{process}{processos} en execució
dins d'un \conceptRef{OS}{sistema operatiu (SO)}. Alguns d'aquests identificadors
s'anomenen \conceptRef{address}{adreces}, precisament perquè s'utilitzen 
per trobar i arribar a un dispositiu remot. 


Les \conceptRef{address}{adreces} solen ser identificadors \textit{numèrics}, és a dir, només un nombre.
Com a tals, es poden expressar en bases diferents, incloent-hi la decimal i l'\concept{hexadecimal}.
Aquests nombres s'extreuen d'un conjunt predefinit, i l'elecció d'aquest conjunt
determina el nombre màxim d'elements que podem identificar de manera única.

\begin{remark}
Les adreces expressades com \otherBase{d8:43:ae:61:ed:f1} o \otherBase{192.168.1.1} 
(com es veurà més endavant al curs) també són nombres que hauríem pogut expressar 
com \otherBase{237785200061937} i \otherBase{3232235777}, respectivament
(cosa que no és tan pràctica).
\end{remark}

\begin{exercise}
\textbf{Quants elements es poden identificar de manera única} (com a màxim) amb una adreça que expressem 
amb 5 bits? I amb 10? I amb 20? Expressa la solució general matemàticament.
\end{exercise}

\begin{exercise}
Suposem que hi ha $1000$~màquines en una LAN. \textbf{Quants bits d'adreça calen}?
I per a $800$~màquines? I per a $1100$? Expressa la solució general matemàticament.
\end{exercise} 

\begin{remark}
El 2019 ens vam quedar sense adreces d'Internet (v4). Ups!
\end{remark}

\section{Com hi arribo?}

En molts escenaris, calen diverses adreces per dur a terme la comunicació 
\conceptRef{end-to-end}{extrem a extrem} desitjada. Per exemple, pot ser que hàgim d'identificar
un únic dispositiu dins d'una \concept{LAN}, i aquesta LAN dins d'una \concept{WAN}. De vegades, caldrà 
identificar un \conceptRef{process}{procés} (un programa en execució) dins d'aquest dispositiu. 
No és gaire diferent de quan algú et ve a visitar en persona:

\begin{minipage}{0.5\linewidth}
\begin{center}
\begin{bytefield}{16}
\bitbox{16}{Joseph Edgar Foreman} \\
\bitbox{16}{123 South Fake St.} \\
\bitbox{16}{Adams County} \\
\bitbox{16}{Ohio}
\end{bytefield}
\end{center}
\end{minipage}
% 
\begin{minipage}{0.5\linewidth}
\begin{center}
\begin{bytefield}{24}
\bitbox{24}{Process \#1717403} \\
\bitbox{24}{Device \#51688798} \\
\bitbox{24}{LAN 3141 (Data Engineering Labs)} \\
\bitbox{24}{WAN 442 (UAB)}
\end{bytefield} 
\end{center}
\end{minipage}

Sigui quin sigui el sistema d'adreçament que fem servir, ha de contenir prou informació per localitzar i arribar 
al destí, \ie, per \conceptRef{routing}{encaminar} els nostres missatges a través del graf de 
connexions. El procés d'encaminament implica diversos passos o ``salts'' entre xarxes
fins que s'arriba a la LAN de destí.

El nom \concept{router} (encaminador) designa un tipus de dispositiu de xarxa
la feina del qual és reenviar aquests missatges entre xarxes. Aquests dispositius han de contenir
prou informació perquè, quan reben un missatge amb una \conceptRef{address}{adreça} de destí,
puguin decidir quina \conceptRef{network interface}{interfície de xarxa} han de fer servir. 
Per la seva banda, els dispositius que no són routers només cal que estiguin configurats per poder arribar al router següent.


\begin{remark}
A casa, el teu \concept{router} normalment també fa el paper de \concept{switch}, però són conceptes diferents.
\end{remark}

% \begin{center}
% \includegraphics[width=0.5\linewidth]{how_many_packets.pdf}
% \end{center}

\begin{exercise}
Aquesta figura representa un grapat de dispositius i les seves connexions físiques.
Aquests dispositius només es poden comunicar enviant missatges a través 
d'aquestes connexions. Tot i així, \node{3} vol enviar un missatge a \node{7}.\\
% 
\begin{center}
\includegraphics[width=0.5\linewidth]{how_many_packets.pdf}\\
\end{center}
% 
\begin{itemize}
 \item Pots identificar algun dispositiu que probablement actuï com a \concept{switch}? I com a \concept{router}?
 \item Quin és el \textbf{nombre mínim de missatges} enviats perquè \node{3} pugui contactar amb \node{7},
 i \node{7} pugui respondre a \node{3}? Enumera'ls en l'ordre en què es produeixen.
 \item Quina informació s'ha d'incloure en aquests missatges perquè la comunicació 
 \node{3} $\leftrightarrow$ \node{7} pugui tenir lloc?
\end{itemize}
\label{ex:how_many_packets}
\end{exercise}
% 

Els \conceptRef{swimlane diagram}{diagrames de carrils} com el de sota s'utilitzen per representar les interaccions
(\eg, intercanvi de missatges) entre actors (\eg, dispositius de xarxa). En aquests diagrames, l'eix vertical representa 
el temps, cosa que és molt valuosa per identificar relacions de causa-efecte.
\begin{center}
\includegraphics[width=0.35\linewidth]{swimlane_example.eps}
\end{center}

\begin{exercise}
A partir de la teva resposta a l'Ex.\ref{ex:how_many_packets}, \textbf{elabora un diagrama de carrils}
que representi els intercanvis de missatges. Per a cada missatge (\eg, damunt de les fletxes), 
inclou-hi la informació d'adreçament que conté.
\end{exercise}


